\documentclass[12pt]{article}
\usepackage[utf8]{inputenc}
\usepackage[T1]{fontenc}
\usepackage{hyperref}
\title{Kenyan Sign Language Language Resources for K-12 Curriculum Development}
\author{Compiled by Claude Sonnet 5.5 (claude-sonnet-5-5) for Andr\'as Kornai}
\date{}
\begin{document}
\maketitle

\section{Introduction}
Kenyan Sign Language (KSL; ISO 639-3 xki; Glottocode keny1241) is the sign language of the Kenyan deaf community. The search results indicate that KSL has only recently been brought into the school curriculum. A curriculum review aimed to make KSL part of the curriculum for all learners, deaf or hearing, and a pilot ran in 470 schools. Wikipedia reports that KSL is not generally used in the classrooms of the residential schools for deaf students, even though it is their main language. That statement may be outdated.

\section{Primary Educational Resources}
\begin{itemize}
\item Global Partnership for Education, on the Kenyan school system trying sign language for all children: \url{https://globalpartnership.org/blog/kenyan-school-system-trying-sign-language-all-children}
\item Deaf Child Worldwide, on sign language for all in Kenyan schools: \url{https://www.ndcs.org.uk/deaf-child-worldwide/our-blog/sign-language-for-all-in-kenyan-schools}
\item People Daily, on a bill to pave the way for sign language use in schools (the bill's status was not clear from the results): \url{https://peopledaily.digital/news/bill-to-pave-way-for-sign-language-use-in-schools/amp}
\item Kenyatta University repository record, a study of KSL in schools (western Kenya primary schools; it recommends KSL-illustrated materials and teacher training): \url{https://ir-library.ku.ac.ke/handle/123456789/7503?show=full}
\end{itemize}

\section{Linguistic and Reference Materials}
\begin{itemize}
\item Wikipedia, Kenyan Sign Language: \url{https://wikipedia.com/wiki/ISO_639:xki}
\item Glottolog reference, Rothenborg-Jensen and Yego 2011: \url{https://glottolog.org/resource/reference/id/320492}
\item K24 article on how KSL works and where to learn it (it contains claims that conflict with Wikipedia, so treat it with caution): \url{https://k24.digital/lifestyle/life-hacks/international-week-of-the-deaf-how-kenyan-sign-language-works-and-where-to-learn-it}
\end{itemize}
Dictionaries mentioned in the search results are a 1991 KSL dictionary, a 2011 English/Kiswahili bilingual dictionary from the Kenya Institute of Special Education, and an online dictionary at ksldictionary.com. The online dictionary was named in a result but not opened, and no separate URL for it appeared in the results.

\section{Religious and Community Resources}
None found in the search results.

\section{Digital and Audio Resources}
A search result mentions an interactive digital dictionary developed by the KSL research project with Peace Corps volunteers. No verified URL for it was found.

\section{Curriculum Development Context}
Policy momentum exists: KSL piloting for all learners and a proposed bill giving the Kenya Institute of Curriculum Development responsibility for materials for learners with hearing disabilities. The gaps appear to be standard KSL-medium textbooks, trained teachers and video materials. These need verifying with Kenyan deaf organisations.
\end{document}